\documentclass[10pt,a4paper]{amsart}
\usepackage[margin=19mm]{geometry}
\usepackage{amsmath,amssymb,mathtools}
\usepackage[scr=rsfs,cal=euler]{mathalfa}
\usepackage{xfrac}
\usepackage[unicode,hidelinks,bookmarks=false]{hyperref}
\newcommand{\ie}{i.e.\ }
\newcommand{\cf}{cf.\ }
\newcommand{\resp}{respectively\ }
\newcommand{\Subsection}{\subsection}
\providecommand{\nobreakdash}{\nobreak\hbox{-}}
\setlength{\emergencystretch}{2em}
\multlinegap0pt
\usepackage{bm}
\usepackage{bbm}
\usepackage{dsfont}
\usepackage{xspace}
\usepackage{cancel}
\usepackage{mathtools}
\usepackage[shortlabels]{enumitem}
\usepackage{amsthm}
\makeatletter
\@ifundefined{theorem}{\newtheorem{theorem}{Theorem}}{}
\@ifundefined{lemma}{\newtheorem{lemma}[theorem]{Lemma}}{}
\makeatother
\newcommand{\Exp}{{\mathbb{E}}}
\newcommand{\Pro}{{\mathbb{P}}}
\title[Efficient Importance Sampling for Wrong Exit Probabilities o]{Efficient Importance Sampling for Wrong Exit Probabilities over Combinatorially Many Rare Regions}
\author{Yanglei Song, Georgios Fellouris}
\date{Source: arXiv:2509.14596v1 (CC BY 4.0)}
\begin{document}
\maketitle

\noindent\emph{Source and licence.} Efficient Importance Sampling for Wrong Exit Probabilities over Combinatorially Many Rare Regions. Version arXiv:2509.14596v1 (CC BY 4.0), distributed under the Creative Commons Attribution 4.0 International licence, as recorded in the arXiv licence metadata for this exact version and shown on the abstract page https://arxiv.org/abs/2509.14596v1. Full rights notice: licenses/arxiv-2509.14596-CC-BY-4.0.txt.\par\smallskip

\noindent\textit{Editorial scope.} Counted unit: the Lemma whose source label is \texttt{aux\_lemma\_overall} in the article (source0.tex lines 2837--2842, source label \texttt{aux\_lemma\_overall}), delivered together with the article's own complete original proof of it (source0.tex lines 2843--2888, one \texttt{proof} environment of 46 lines). No mathematics is written, repaired or re-derived: statement and proof are the article's own text line for line. Required context retained uncounted: assumptions\_directions\_of\_mean (lines 631--635); assumption\_on\_X (lines 638--644); assumption\_on\_regions (lines 669--674); assumption:negative\_results (lines 959--964); lemma:aux\_prob1 (lines 2659--2666); lemma:aux\_prob2 (lines 2757--2763); lemma:aux\_prob34 (lines 2778--2791). Every definition, display and label of those lines is retained unchanged. Omitted from this package and not redistributed: 1--630, 636--637, 645--668, 675--958, 965--2658, 2667--2756, 2764--2777, 2792--2836, 2889--3549. Nothing inside a retained excerpt is omitted or altered: every retained body is delivered exactly as the article's lines are written, so no omission receipt of any kind accompanies this package. Citations: the retained excerpts carry no citation command at all, so no bibliography entry of the article is transcribed and no reference is redistributed. Reference adaptation: none was needed and none was made; every reference inside the delivered unit resolves inside it, so no unresolved reference or citation is produced. Printed slips kept literal: no printed word, symbol, hypothesis or number of the counted statement or of its proof was altered. Layout and class: the article's own class is \texttt{imsart} with packages including mathtools, amsthm, amsmath, amsfonts, amssymb, natbib, hyperref, graphicx, enumitem, comment, bbm; the wrapper is the lane's pinned \texttt{amsart} editorial wrapper at 10pt on A4, which restates the theorem-like environments the retained text needs and adds the article's own macro definitions that the retained text needs (\texttt{\textbackslash Exp}, \texttt{\textbackslash Pro}), with extra wrapper packages (bm, bbm, dsfont, xspace, cancel, mathtools, enumitem). The delivered numbering is the unit's own. Assets: no retained span contains a figure environment, a graphics-inclusion command, a PDF-page inclusion, a table or a TikZ picture; no image file is copied into this package, the package declares no asset dependency and no image file is redistributed with it. Licence: version arXiv:2509.14596v1 of this article is distributed under the Creative Commons Attribution 4.0 International licence, as recorded in the arXiv licence metadata for this exact version and shown on the abstract page https://arxiv.org/abs/2509.14596v1. A full rights notice is delivered at licenses/arxiv-2509.14596-CC-BY-4.0.txt. Characters: every retained excerpt is pure ASCII, so the delivered engine, XeLaTeX with its default Latin Modern text font, reports no missing character.
\begin{align}\tag{C1}\label{assumptions_directions_of_mean}
\begin{split}
\text{there exists } \;\; \delta>0 \;\; \text{such that} \;\;  \operatorname{cone}(B(\Exp[X_1], \delta)) \cap W^j = \emptyset \;\; \text{for every } j \in [J].
\end{split}
\end{align}

\begin{align} \tag{C2} \label{assumption_on_X}
    \begin{split}
        &\operatorname{dom}(\Lambda) \text{ is an open set}; \\
        %\text{ and }
         &\operatorname{int}\left(\operatorname{cone}\left( \operatorname{dom}(\Lambda^*)\right)\right) \cap W^j \neq \emptyset, \quad \text{for  every }  j \in [J].
    \end{split}
\end{align}

\begin{align} \tag{C3} \label{assumption_on_regions}
    \begin{split}
        &\operatorname{cone}(W^j) \cap \operatorname{cone}(W^{j'}) = \emptyset \quad \text{for all } 0 \leq j \neq j' \leq J;\\
        &\text{there exists } \;\; \delta>0 \;\; \text{such that} \;\; B(\mathbf{0}, \delta) \cap W^j = \emptyset \quad \text{for each } j \in \{0\} \cup [J].
    \end{split}
\end{align}

\begin{align}\label{assumption:negative_results}  \tag{C4}
    \begin{split}
    &  \overline{\operatorname{cone}(W^j)} \;\cap\; \overline{\operatorname{cone}(W^{j'})} = \{\mathbf{0}\}, \;\;\; \text{ for } \;0 \leq j \neq j' \leq J. \\
    & \text{ Either } \; (i). \; W^0 = \emptyset, \quad  \text{ or } \;\; (ii).\; W^0 \text{ is convex and } \operatorname{int}\left(\operatorname{cone}\left( \operatorname{dom}(\Lambda^*)\right)\right) \cap W^0 \neq \emptyset.
    \end{split}
\end{align}

\begin{lemma}\label{lemma:aux_prob1}
    Fix $j \in [J]$. 
    Assume that conditions \eqref{assumptions_directions_of_mean}, \eqref{assumption_on_X}, and \eqref{assumption_on_regions} hold.  Then 
    for any $\delta > 0$,
        \begin{align*}
              \lim_{b \to \infty} \frac{1}{b}  \log \Pro_{\beta^j}\left(\left|{\tau_b^j}/{b} - \rho^{j}\right| \leq \delta  \right) = 0.
        \end{align*}
\end{lemma}

\begin{lemma}\label{lemma:aux_prob2}
    Fix $j \in [J]$. 
    Assume that conditions \eqref{assumptions_directions_of_mean}, \eqref{assumption_on_X},  \eqref{assumption_on_regions}, and \eqref{assumption:negative_results} hold.  For $j' \in \{0,1,\ldots,J\}$ such that $j' \neq j$,
        \begin{align*}
          \limsup_{b \to \infty} \frac{1}{b}  \log \Pro_{\beta^j}\left(\tau_b^{j'} < \infty\right) < 0.
        \end{align*}  
\end{lemma}

\begin{lemma}\label{lemma:aux_prob34}
    Fix $j \in [J]$. 
    Assume that conditions \eqref{assumptions_directions_of_mean}, \eqref{assumption_on_X}, and \eqref{assumption_on_regions} hold.  
    \begin{enumerate}[label=(\roman*)]
        \item For any $\delta >0$,
        \begin{align*}
          \limsup_{n \to \infty} \frac{1}{n}  \log \Pro_{\beta^j}\left( \frac{1}{n}\beta^j \cdot S_n -\frac{r_j}{\rho^j} \geq \delta  \right) < 0.
        \end{align*}
        \item For any $\theta \in \operatorname{dom}(\Lambda)$ such that $\theta \neq \beta^j$, there exists a sufficiently small $\delta > 0$ such that
        \begin{align*}
          \limsup_{n \to \infty} \frac{1}{n}  \log \Pro_{\beta^j}\left( \frac{1}{n}(\theta -\beta^j) \cdot S_n -\Lambda(\theta) \geq -\delta  \right) < 0.
        \end{align*}      
    \end{enumerate}
\end{lemma}

\begin{lemma}\label{aux_lemma_overall}
    Assume that conditions \eqref{assumptions_directions_of_mean}, \eqref{assumption_on_X},  \eqref{assumption_on_regions}, and \eqref{assumption:negative_results} hold.    Let   $\Theta \subset \operatorname{dom}(\Lambda)$ be a finite set such that $\beta^j \not \in \Theta$. There exist  $\epsilon, \delta > 0$ such that
    \begin{align*}
       \lim_{b \to \infty}  \frac{1}{b} \log \Pro_{\beta^j}\left(\mathcal{A}_{b,\epsilon,\delta,\Theta}^{(j,*)} \right) = 0.
    \end{align*}
\end{lemma}

\begin{proof}
By definition and the union bound,
\begin{align*}
\Pro_{\beta^j}\left(\mathcal{A}_{b,\epsilon,\delta,\Theta}^{(j,*)} \right) \geq &\Pro_{\beta^j}\left(\mathcal{A}_{b,\epsilon}^{(j,2)}\right) - 
\Pro\left( \left(\mathcal{A}_{b}^{(j,1)}\right)^c \right) \\
&- \Pro\left( \mathcal{A}_{b,\epsilon}^{(j,2)}  \setminus \mathcal{A}_{b,\delta}^{(j,3)} \right) -
\Pro\left( \mathcal{A}_{b,\epsilon}^{(j,2)} \setminus \mathcal{A}_{b,\delta,\Theta}^{(j,4)} \right).
\end{align*}
Thus, it suffices to show that there exist $\epsilon,\delta > 0$ such that
\begin{align}
\begin{split}
     &\liminf_{b \to \infty}  \frac{1}{b} \log    \Pro_{\beta^j}\left(\left|{\tau_b^j}/{b} - \rho^{j}\right| \leq \epsilon  \right) = 0, \\
 & \limsup_{b \to \infty}  \frac{1}{b} \log    \Pro_{\beta^j}\left(\tau_b^{k} < \infty \right) < 0  \text{ for each } k \neq j, \\
  & \limsup_{b \to \infty}  \frac{1}{b} \log    \Pro_{\beta^j}\left(\left|{\tau_b^j}/{b} - \rho^{j}\right| \leq   \epsilon , \;\; \beta^j \cdot S_{\tau_b^j} \geq  b(r_j + \delta/4) \right) < 0, \\
   & \limsup_{b \to \infty}  \frac{1}{b} \log    \Pro_{\beta^j}\left(\left|{\tau_b^j}/{b} - \rho^{j}\right| \leq \epsilon,\; (\theta-\beta^j) \cdot S_{\tau_b^j} - \tau_b^j \Lambda(\theta)  \geq 
-\delta b  \right) < 0, \text{ for } \theta \in \Theta.
\end{split}
\label{eqref:aux_overall}
\end{align}
The first two claims have been established in Lemma \ref{lemma:aux_prob1} and Lemma \ref{lemma:aux_prob2}, respectively, for any $\epsilon > 0$.

Now, we focus on the third claim in \eqref{eqref:aux_overall}. By definition and the union bound,
\begin{align*}
 \Pro_{\beta^j}\left(\left|{\tau_b^j}/{b} - \rho^{j}\right| \leq   \epsilon, \;\; \beta^j \cdot S_{\tau_b^j} \geq  b(r_j + \delta/4) \right) 
 \leq \sum_{n = b(\rho^j-\epsilon)}^{b(\rho^j+\epsilon)}  \Pro_{\beta^j}\left(\beta^j \cdot S_{n} \geq  b(r_j + \delta/4) \right).
\end{align*}
For integer $n \geq 0$ such that $|n/b - \rho^j| \leq \epsilon$, we have
\begin{align*}
    \frac{1}{n}  b(r_j + \epsilon/4) \geq \frac{1}{\rho^j +\epsilon} (r_j + \delta/4) = \frac{r_j}{\rho^j} + \tilde{\delta},
\end{align*}
where we define
$$
\tilde{\delta}:= \frac{1}{\rho^j + \epsilon}\left[\frac{\delta}{4} - \frac{r_j}{\rho^j} \epsilon \right].
$$
Thus, if $\delta/4 >  ({r_j}/{\rho^j}) \epsilon $, we have $\tilde{\delta} > 0$, and by Lemma \ref{lemma:aux_prob34}(i), there exists some constant $c > 0$ such that  for all sufficiently large $n$, 
\begin{align*}
           \frac{1}{n}  \log \Pro_{\beta^j}\left( \frac{1}{n}\beta^j \cdot S_n \geq \frac{r_j}{\rho^j} + \tilde{\delta}  \right) \leq -c < 0.
    \end{align*}
As a result, for sufficiently large $b > 0$, if $\delta/4 >  ({r_j}/{\rho^j}) \epsilon $,
\begin{align*}
    \sum_{n = b(\rho^j-\epsilon)}^{b(\rho^j+\epsilon)}  \Pro_{\beta^j}\left(\beta^j \cdot S_{n} \geq  b(r_j + \delta/4) \right) \leq (2b\epsilon) e^{-c b(\rho^j - \epsilon)}.
\end{align*}
Thus, if $\delta/4 >  ({r_j}/{\rho^j}) \epsilon $ and $\epsilon < \rho^j$, then the third claim in \eqref{eqref:aux_overall} holds.

The fourth claim in \eqref{eqref:aux_overall} follows from similar arguments as above, once we replace   Lemma \ref{lemma:aux_prob34}(i) with Lemma \ref{lemma:aux_prob34}(ii).
\end{proof}

\end{document}
